% Header: General study / work / project template
%%%%%%%%%%%%%%%%%%%%%%%%%%%%%%%%%%%%%%%%%%%%%%%%%%%%%%%%%%%%%%%%%%%

\documentclass[11pt,a4paper]{scrartcl}

% Layout
\usepackage[margin=2.5cm]{geometry}
\usepackage[onehalfspacing]{setspace}

% General packages
\usepackage{graphicx}
\usepackage{enumitem}
\usepackage{tcolorbox}
\tcbuselibrary{breakable}
\usepackage[breaklinks=true,colorlinks=true,linkcolor=blue,urlcolor=blue,citecolor=blue]{hyperref}
\usepackage{amsmath,amsthm,amssymb,mathtools}
\usepackage{icomma}
\usepackage[T1]{fontenc}
\usepackage[utf8]{inputenc}
\usepackage[ngerman]{babel}

% Units / tables / floats / references
\usepackage[locale = DE,space-before-unit=true,per-mode = symbol]{siunitx}
\usepackage{booktabs,multirow}
\usepackage{placeins}
\usepackage[natbib,abbreviate=true,doi=false,style=numeric-comp,giveninits=true,sorting=none]{biblatex}
\usepackage{csquotes}
\usepackage{pgfplots}
\pgfplotsset{compat=1.18}

% Optional bibliography
% \addbibresource{references.bib}

% Graphics paths
\graphicspath{{Bilder/} {images/} {figures/} {chapters/}}

% Optional math shortcuts
\newcommand{\R}{\mathbb{R}}
\newcommand{\N}{\mathbb{N}}
\newcommand{\Q}{\mathbb{Q}}
\newcommand{\set}[1]{\left\{ #1 \right\}}
\newcommand{\abs}[1]{\left| #1 \right|}
\newcommand{\norm}[1]{\left\lVert #1 \right\rVert}
\newcommand{\ol}[1]{\overline{#1}}
\newcommand{\half}{\frac{1}{2}}

% Optional units
\DeclareSIUnit{\year}{a}
\DeclareSIUnit{\day}{d}

% Box styles
\newtcolorbox{calculationbox}[1][]{
    title=Calculation,
    breakable,
    #1
}

\newtcolorbox{solutionbox}[1][]{
    title=Solution,
    breakable,
    #1
}

\newtcolorbox{answerbox}[1][]{
    title=Answer,
    breakable,
    #1
}

\newtcolorbox{notebox}[1][]{
    title=Notes,
    breakable,
    #1
}

\newtcolorbox{sidecalcbox}[1][]{
    title=Side Calculation,
    breakable,
    #1
}

%%%%%%%%%%%%%%%%%%%%%%%%%%%%%%%%%%%%%%%%%%%%%%%%%%%%%%%%%%%%%%%%%%%
\begin{document}

    \title{Physik 4 - Blatt 07}
    \author{Tobias Laser}
    \date{\today}
    \maketitle
    \vfill
    \thispagestyle{empty}
    \newpage

    \pagenumbering{arabic}

    \paragraph{Aufgabe 1: Radioaktivität}

    \begin{enumerate}[label=(\alph*)]
        \item
            In seiner in der Natur vorkommenden Zusammensetzung besteht Uran zu
            \SI{99.28}{\percent} aus ${}^{238}\mathrm{U}$ mit
            $t_{1/2}({}^{238}\mathrm{U})=4{,}468\cdot 10^9\,\mathrm{a}$ und zu
            \SI{0.72}{\percent} aus ${}^{235}\mathrm{U}$ mit
            $t_{1/2}({}^{235}\mathrm{U})=7{,}038\cdot 10^8\,\mathrm{a}$.
            Wie alt musste die Materie des Sonnensystems sein, wenn man annimmt,
            dass bei seiner Entstehung beide Isotope in gleicher Häufigkeit vorhanden waren?
            Wie interpretieren Sie das Ergebnis?

            \begin{calculationbox}
                Wir verwenden das radioaktive Zerfallsgesetz
                \[
                    N(t)=N(0)e^{-\lambda t}
                \]
                mit der Zerfallskonstante
                \[
                    \lambda=\frac{\ln 2}{t_{1/2}}.
                \]
                Zur Entstehung seien beide Isotope gleich häufig gewesen:
                \[
                    N_{238}(0)=N_{235}(0)=N_0.
                \]
                Heute gilt daher
                \begin{align*}
                    N_{238}(t)&=N_0 e^{-\lambda_{238}t},\\
                    N_{235}(t)&=N_0 e^{-\lambda_{235}t}.
                \end{align*}
                Wir bilden den Quotienten, damit sich $N_0$ herauskürzt:
                \begin{align*}
                    \frac{N_{238}(t)}{N_{235}(t)}
                    &=\frac{N_0e^{-\lambda_{238}t}}{N_0e^{-\lambda_{235}t}}\\
                    &=e^{(\lambda_{235}-\lambda_{238})t}.
                \end{align*}
                Der heutige Quotient ist durch die natürliche Zusammensetzung gegeben:
                \[
                    \frac{N_{238}(t)}{N_{235}(t)}=\frac{99{,}28}{0{,}72}=137{,}89.
                \]
                Also folgt
                \begin{align*}
                    137{,}89&=e^{(\lambda_{235}-\lambda_{238})t},\\
                    \ln(137{,}89)&=(\lambda_{235}-\lambda_{238})t,\\
                    t&=\frac{\ln(137{,}89)}{\lambda_{235}-\lambda_{238}}.
                \end{align*}
                Nun berechnen wir die Zerfallskonstanten:
                \begin{align*}
                    \lambda_{238}
                    &=\frac{\ln 2}{4{,}468\cdot 10^9\,\mathrm{a}}
                    =1{,}551\cdot 10^{-10}\,\mathrm{a}^{-1},\\
                    \lambda_{235}
                    &=\frac{\ln 2}{7{,}038\cdot 10^8\,\mathrm{a}}
                    =9{,}849\cdot 10^{-10}\,\mathrm{a}^{-1}.
                \end{align*}
                Einsetzen liefert
                \begin{align*}
                    t
                    &=\frac{\ln(137{,}89)}{9{,}849\cdot10^{-10}-1{,}551\cdot10^{-10}}\,\mathrm{a}\\
                    &=5{,}94\cdot 10^9\,\mathrm{a}.
                \end{align*}
            \end{calculationbox}

            \begin{answerbox}
                Unter der Annahme gleicher Anfangshäufigkeit wäre die Materie des Sonnensystems etwa
                \[
                    \boxed{t\approx 5{,}94\cdot 10^9\,\mathrm{Jahre}}
                \]
                alt.

                Interpretation: Dieses Ergebnis ist größer als das übliche Alter des Sonnensystems.
                Die Annahme, dass ${}^{238}\mathrm{U}$ und ${}^{235}\mathrm{U}$ bei der Entstehung exakt gleich häufig vorhanden waren,
                ist also eine starke Modellannahme. Das Ergebnis zeigt vor allem, dass ${}^{235}\mathrm{U}$ wegen seiner kürzeren Halbwertszeit deutlich schneller zerfällt als ${}^{238}\mathrm{U}$.
            \end{answerbox}

        \item
            Am Tatort eines Verbrechens wird ein verschlossener Behälter gefunden,
            in dem sich ${}^{210}\mathrm{Po}$ und \SI{29.62}{\micro\gram} von dessen Tochterelement befinden.
            Der Inhalt des Behälters hat die Aktivität $A=1{,}56\cdot 10^8\,\mathrm{Bq}$.
            Schreiben Sie die Zerfallsgleichung für den Alphastrahler ${}^{210}\mathrm{Po}$ hin und berechnen Sie,
            vor wie vielen Tagen das Polonium hergestellt wurde. Die Halbwertszeit beträgt \SI{138.38}{d}.

            \begin{calculationbox}
                Beim Alpha-Zerfall wird ein Heliumkern ${}^{4}_{2}\mathrm{He}$ ausgesendet.
                Dadurch sinkt die Massenzahl um $4$ und die Ordnungszahl um $2$.
                Für Polonium gilt $Z=84$, also entsteht Blei mit $Z=82$:
                \[
                    \boxed{
                    {}^{210}_{84}\mathrm{Po}\longrightarrow {}^{206}_{82}\mathrm{Pb}+{}^{4}_{2}\mathrm{He}
                    }
                \]
            \end{calculationbox}

            \begin{calculationbox}
                Die Aktivität des noch vorhandenen Poloniums ist
                \[
                    A(t)=\lambda N_{\mathrm{Po}}(t).
                \]
                Damit ist
                \[
                    N_{\mathrm{Po}}(t)=\frac{A(t)}{\lambda}.
                \]
                Die Zerfallskonstante ist
                \begin{align*}
                    \lambda
                    &=\frac{\ln 2}{t_{1/2}}
                    =\frac{\ln 2}{138{,}38\cdot 24\cdot 3600\,\mathrm{s}}\\
                    &=5{,}798\cdot10^{-8}\,\mathrm{s}^{-1}.
                \end{align*}
                Also sind aktuell noch
                \begin{align*}
                    N_{\mathrm{Po}}(t)
                    &=\frac{1{,}56\cdot10^8\,\mathrm{s}^{-1}}{5{,}798\cdot10^{-8}\,\mathrm{s}^{-1}}\\
                    &=2{,}69\cdot10^{15}
                \end{align*}
                Poloniumkerne vorhanden.
            \end{calculationbox}

            \begin{calculationbox}
                Das Tochterelement ist ${}^{206}\mathrm{Pb}$. Aus der Masse des Bleis berechnen wir die Anzahl der Bleikerne:
                \[
                    N_{\mathrm{Pb}}=\frac{m_{\mathrm{Pb}}}{M_{\mathrm{Pb}}}N_A.
                \]
                Mit
                \[
                    m_{\mathrm{Pb}}=29{,}62\cdot10^{-6}\,\mathrm{g},
                    \qquad
                    M_{\mathrm{Pb}}\approx 206\,\mathrm{g/mol}
                \]
                folgt
                \begin{align*}
                    N_{\mathrm{Pb}}
                    &=\frac{29{,}62\cdot10^{-6}\,\mathrm{g}}{206\,\mathrm{g/mol}}
                    \cdot 6{,}022\cdot10^{23}\,\mathrm{mol}^{-1}\\
                    &=8{,}66\cdot10^{16}.
                \end{align*}
            \end{calculationbox}

            \begin{calculationbox}
                Da der Behälter verschlossen war und wir annehmen, dass am Anfang nur Polonium vorhanden war, gilt:
                \[
                    N_0=N_{\mathrm{Po}}(t)+N_{\mathrm{Pb}}(t).
                \]
                Außerdem gilt
                \[
                    N_{\mathrm{Po}}(t)=N_0e^{-\lambda t}.
                \]
                Damit
                \begin{align*}
                    \frac{N_{\mathrm{Po}}(t)}{N_0}&=e^{-\lambda t},\\
                    \frac{N_{\mathrm{Po}}(t)}{N_{\mathrm{Po}}(t)+N_{\mathrm{Pb}}(t)}&=e^{-\lambda t}.
                \end{align*}
                Logarithmieren und nach $t$ auflösen:
                \begin{align*}
                    t
                    &=-\frac{1}{\lambda}
                    \ln\left(
                    \frac{N_{\mathrm{Po}}(t)}{N_{\mathrm{Po}}(t)+N_{\mathrm{Pb}}(t)}
                    \right)\\
                    &=-\frac{1}{5{,}798\cdot10^{-8}\,\mathrm{s}^{-1}}
                    \ln\left(
                    \frac{2{,}69\cdot10^{15}}{2{,}69\cdot10^{15}+8{,}66\cdot10^{16}}
                    \right)\\
                    &=6{,}04\cdot10^7\,\mathrm{s}\\
                    &=699\,\mathrm{d}.
                \end{align*}
            \end{calculationbox}

            \begin{answerbox}
                Das Polonium wurde vor ungefähr
                \[
                    \boxed{699\,\mathrm{Tagen}}
                \]
                hergestellt. Das entspricht ungefähr \SI{1.91}{Jahren}.
            \end{answerbox}
    \end{enumerate}

    \newpage
    \paragraph{Aufgabe 2: Kettenzerfälle -- Theorie}

    Ein Radionuklid $X_1$ mit Zerfallskonstante $\lambda_1$ zerfällt in die Tochtersubstanz $X_2$.
    Diese zerfällt mit Zerfallskonstante $\lambda_2$ weiter in das stabile Nuklid $X_3$.

    \begin{enumerate}[label=(\alph*)]
        \item
            Formulieren Sie die Differentialgleichungen für die Anzahl $N_1(t)$ und $N_2(t)$ der Kerne $X_1$ und $X_2$.
            Lösen Sie beide Differentialgleichungen unter der Annahme, dass zur Zeit $t=0$ nur $N_0$ Kerne $X_1$ vorhanden sind.

            \begin{calculationbox}
                Die Zerfallskette lautet
                \[
                    X_1 \xrightarrow{\lambda_1} X_2 \xrightarrow{\lambda_2} X_3.
                \]
                Für $X_1$ gibt es nur Zerfall, aber keine Produktion. Daher gilt
                \[
                    \frac{dN_1}{dt}=-\lambda_1N_1.
                \]
                Für $X_2$ gibt es zwei Beiträge:
                \begin{itemize}
                    \item Produktion durch Zerfall von $X_1$: $+\lambda_1N_1$,
                    \item Verlust durch Zerfall von $X_2$: $-\lambda_2N_2$.
                \end{itemize}
                Also gilt
                \[
                    \frac{dN_2}{dt}=\lambda_1N_1-\lambda_2N_2.
                \]
                Die Anfangsbedingungen sind
                \[
                    N_1(0)=N_0,
                    \qquad
                    N_2(0)=0,
                    \qquad
                    N_3(0)=0.
                \]
            \end{calculationbox}

            \begin{calculationbox}
                Zuerst lösen wir die Gleichung für $N_1$:
                \begin{align*}
                    \frac{dN_1}{dt}&=-\lambda_1N_1,\\
                    \frac{dN_1}{N_1}&=-\lambda_1dt.
                \end{align*}
                Integration liefert
                \[
                    \ln N_1=-\lambda_1t+C.
                \]
                Exponentieren ergibt
                \[
                    N_1(t)=Ce^{-\lambda_1t}.
                \]
                Mit $N_1(0)=N_0$ folgt $C=N_0$. Also
                \[
                    \boxed{N_1(t)=N_0e^{-\lambda_1t}.}
                \]
            \end{calculationbox}

            \begin{calculationbox}
                Jetzt setzen wir $N_1(t)$ in die Differentialgleichung für $N_2$ ein:
                \[
                    \frac{dN_2}{dt}=\lambda_1N_0e^{-\lambda_1t}-\lambda_2N_2.
                \]
                Wir bringen die Gleichung in Standardform:
                \[
                    \frac{dN_2}{dt}+\lambda_2N_2=\lambda_1N_0e^{-\lambda_1t}.
                \]
                Das ist eine lineare Differentialgleichung erster Ordnung.
                Der integrierende Faktor ist
                \[
                    \mu(t)=e^{\lambda_2t}.
                \]
                Multiplikation mit $e^{\lambda_2t}$ ergibt
                \begin{align*}
                    e^{\lambda_2t}\frac{dN_2}{dt}+\lambda_2e^{\lambda_2t}N_2
                    &=\lambda_1N_0e^{(\lambda_2-\lambda_1)t},\\
                    \frac{d}{dt}\left(e^{\lambda_2t}N_2(t)\right)
                    &=\lambda_1N_0e^{(\lambda_2-\lambda_1)t}.
                \end{align*}
                Nun integrieren wir von $0$ bis $t$:
                \begin{align*}
                    e^{\lambda_2t}N_2(t)-N_2(0)
                    &=\lambda_1N_0\int_0^t e^{(\lambda_2-\lambda_1)t'}\,dt'.
                \end{align*}
                Wegen $N_2(0)=0$ folgt für $\lambda_1\neq\lambda_2$:
                \begin{align*}
                    e^{\lambda_2t}N_2(t)
                    &=\lambda_1N_0
                    \left[\frac{e^{(\lambda_2-\lambda_1)t'} }{\lambda_2-\lambda_1}\right]_0^t\\
                    &=\frac{\lambda_1N_0}{\lambda_2-\lambda_1}
                    \left(e^{(\lambda_2-\lambda_1)t}-1\right).
                \end{align*}
                Multiplikation mit $e^{-\lambda_2t}$ liefert
                \begin{align*}
                    N_2(t)
                    &=\frac{\lambda_1N_0}{\lambda_2-\lambda_1}
                    \left(e^{-\lambda_1t}-e^{-\lambda_2t}\right).
                \end{align*}
            \end{calculationbox}

            \begin{answerbox}
                Für $\lambda_1\neq\lambda_2$ gilt
                \[
                    \boxed{N_1(t)=N_0e^{-\lambda_1t}}
                \]
                und
                \[
                    \boxed{
                    N_2(t)=\frac{\lambda_1N_0}{\lambda_2-\lambda_1}
                    \left(e^{-\lambda_1t}-e^{-\lambda_2t}\right).
                    }
                \]
            \end{answerbox}

        \item
            Bestimmen Sie nun die Anzahl $N_3(t)$.

            \begin{calculationbox}
                Da $X_3$ stabil ist, verschwinden keine $X_3$-Kerne mehr.
                Außerdem geht jeder ursprüngliche $X_1$-Kern entweder noch als $X_1$, als $X_2$ oder schon als $X_3$ vor.
                Daher gilt die Erhaltung der Gesamtzahl der Kerne:
                \[
                    N_1(t)+N_2(t)+N_3(t)=N_0.
                \]
                Also
                \[
                    N_3(t)=N_0-N_1(t)-N_2(t).
                \]
                Einsetzen der Ergebnisse aus Teil (a) liefert
                \begin{align*}
                    N_3(t)
                    &=N_0-N_0e^{-\lambda_1t}
                    -\frac{\lambda_1N_0}{\lambda_2-\lambda_1}
                    \left(e^{-\lambda_1t}-e^{-\lambda_2t}\right).
                \end{align*}
            \end{calculationbox}

            \begin{answerbox}
                Für $\lambda_1\neq\lambda_2$ ist
                \[
                    \boxed{
                    N_3(t)=N_0-N_0e^{-\lambda_1t}
                    -\frac{\lambda_1N_0}{\lambda_2-\lambda_1}
                    \left(e^{-\lambda_1t}-e^{-\lambda_2t}\right).
                    }
                \]
                Äquivalent kann man schreiben:
                \[
                    \boxed{
                    N_3(t)=N_0\left[1+
                    \frac{\lambda_1e^{-\lambda_2t}-\lambda_2e^{-\lambda_1t}}
                    {\lambda_2-\lambda_1}
                    \right].
                    }
                \]
            \end{answerbox}

        \item
            Was folgt für $N_2(t)$ im Sonderfall, dass $\lambda_1=\lambda_2$ ist?

            \begin{calculationbox}
                Ausgangspunkt ist
                \[
                    N_2(t)=\frac{\lambda_1N_0}{\lambda_2-\lambda_1}
                    \left(e^{-\lambda_1t}-e^{-\lambda_2t}\right).
                \]
                Wir setzen
                \[
                    \lambda_1=\lambda,
                    \qquad
                    \lambda_2\to\lambda.
                \]
                Dann erhalten wir formal den Ausdruck
                \[
                    N_2(t)=\lambda N_0
                    \frac{e^{-\lambda t}-e^{-\lambda_2t}}{\lambda_2-
                    \lambda}.
                \]
                Für $\lambda_2\to\lambda$ entsteht ein Ausdruck vom Typ $0/0$.
                Daher verwenden wir die Regel von l'Hospital:
                \begin{align*}
                    \lim_{\lambda_2\to\lambda}
                    \frac{e^{-\lambda t}-e^{-\lambda_2t}}{\lambda_2-
                    \lambda}
                    &=\lim_{\lambda_2\to\lambda}
                    \frac{\frac{d}{d\lambda_2}\left(e^{-\lambda t}-e^{-\lambda_2t}\right)}
                    {\frac{d}{d\lambda_2}\left(\lambda_2-\lambda\right)}\\
                    &=\lim_{\lambda_2\to\lambda}
                    \frac{t e^{-\lambda_2t}}{1}\\
                    &=t e^{-\lambda t}.
                \end{align*}
                Damit folgt
                \[
                    N_2(t)=\lambda N_0 t e^{-\lambda t}.
                \]
            \end{calculationbox}

            \begin{answerbox}
                Im Sonderfall $\lambda_1=\lambda_2=\lambda$ gilt
                \[
                    \boxed{N_2(t)=N_0\lambda t e^{-\lambda t}.}
                \]
            \end{answerbox}

        \item
            Zeichnen Sie Diagramme der Aktivitäten in Einheiten von $\lambda_1N_0$ für die beiden Fälle
            $\lambda_2=\lambda_1/5$ und $\lambda_2=5\lambda_1$ im Bereich $\lambda_1t=0\ldots5$.
            Stellen Sie diese Diagramme sowohl mit logarithmischer als auch linearer $y$-Achse dar.

            \begin{calculationbox}
                Die Aktivität eines Radionuklids ist allgemein
                \[
                    A_i(t)=\lambda_iN_i(t).
                \]
                Wir führen die dimensionslose Variable
                \[
                    x=\lambda_1t
                \]
                und das Verhältnis
                \[
                    \mu=\frac{\lambda_2}{\lambda_1}
                \]
                ein. Dann gilt
                \[
                    \frac{A_1}{\lambda_1N_0}=e^{-x}.
                \]
                Für $A_2$ gilt
                \begin{align*}
                    \frac{A_2}{\lambda_1N_0}
                    &=\frac{\lambda_2N_2}{\lambda_1N_0}\\
                    &=\frac{\lambda_2}{\lambda_1N_0}
                    \frac{\lambda_1N_0}{\lambda_2-\lambda_1}
                    \left(e^{-\lambda_1t}-e^{-\lambda_2t}\right)\\
                    &=\frac{\mu}{\mu-1}\left(e^{-x}-e^{-\mu x}\right).
                \end{align*}
                Das stabile Nuklid $X_3$ hat keine Aktivität, weil es nicht weiter zerfällt:
                \[
                    A_3=0.
                \]
            \end{calculationbox}

            \begin{answerbox}
                Die zu zeichnenden dimensionslosen Funktionen sind
                \[
                    \boxed{\frac{A_1}{\lambda_1N_0}=e^{-x}}
                \]
                und
                \[
                    \boxed{\frac{A_2}{\lambda_1N_0}=\frac{\mu}{\mu-1}\left(e^{-x}-e^{-\mu x}\right),
                    \qquad x=\lambda_1t,
                    \qquad \mu=\frac{\lambda_2}{\lambda_1}.}
                \]
                Für die zwei Fälle ist also
                \[
                    \mu=\frac15
                    \qquad\text{oder}\qquad
                    \mu=5.
                \]
            \end{answerbox}

            \begin{figure}[h]
                \centering
                \begin{tikzpicture}
                    \begin{axis}[
                        width=0.88\textwidth,
                        height=0.48\textwidth,
                        xmin=0,xmax=5,
                        ymin=0,ymax=1.05,
                        xlabel={$x=\lambda_1t$},
                        ylabel={$A/(\lambda_1N_0)$},
                        legend pos=north east,
                        grid=both,
                        title={Linear: $\lambda_2=\lambda_1/5$}
                    ]
                    \addplot[domain=0:5,samples=200,thick] {exp(-x)};
                    \addlegendentry{$A_1$}
                    \addplot[domain=0:5,samples=200,thick,dashed] {(0.2/(0.2-1))*(exp(-x)-exp(-0.2*x))};
                    \addlegendentry{$A_2$}
                    \end{axis}
                \end{tikzpicture}
                \caption{Aktivitäten für $\lambda_2=\lambda_1/5$ mit linearer $y$-Achse.}
            \end{figure}

            \begin{figure}[h]
                \centering
                \begin{tikzpicture}
                    \begin{semilogyaxis}[
                        width=0.88\textwidth,
                        height=0.48\textwidth,
                        xmin=0,xmax=5,
                        ymin=1e-3,ymax=1.05,
                        xlabel={$x=\lambda_1t$},
                        ylabel={$A/(\lambda_1N_0)$},
                        legend pos=north east,
                        grid=both,
                        title={Logarithmisch: $\lambda_2=\lambda_1/5$}
                    ]
                    \addplot[domain=0:5,samples=200,thick] {exp(-x)};
                    \addlegendentry{$A_1$}
                    \addplot[domain=0.001:5,samples=200,thick,dashed] {(0.2/(0.2-1))*(exp(-x)-exp(-0.2*x))};
                    \addlegendentry{$A_2$}
                    \end{semilogyaxis}
                \end{tikzpicture}
                \caption{Aktivitäten für $\lambda_2=\lambda_1/5$ mit logarithmischer $y$-Achse.}
            \end{figure}

            \begin{figure}[h]
                \centering
                \begin{tikzpicture}
                    \begin{axis}[
                        width=0.88\textwidth,
                        height=0.48\textwidth,
                        xmin=0,xmax=5,
                        ymin=0,ymax=1.05,
                        xlabel={$x=\lambda_1t$},
                        ylabel={$A/(\lambda_1N_0)$},
                        legend pos=north east,
                        grid=both,
                        title={Linear: $\lambda_2=5\lambda_1$}
                    ]
                    \addplot[domain=0:5,samples=200,thick] {exp(-x)};
                    \addlegendentry{$A_1$}
                    \addplot[domain=0:5,samples=200,thick,dashed] {(5/(5-1))*(exp(-x)-exp(-5*x))};
                    \addlegendentry{$A_2$}
                    \end{axis}
                \end{tikzpicture}
                \caption{Aktivitäten für $\lambda_2=5\lambda_1$ mit linearer $y$-Achse.}
            \end{figure}

            \begin{figure}[h]
                \centering
                \begin{tikzpicture}
                    \begin{semilogyaxis}[
                        width=0.88\textwidth,
                        height=0.48\textwidth,
                        xmin=0,xmax=5,
                        ymin=1e-3,ymax=1.05,
                        xlabel={$x=\lambda_1t$},
                        ylabel={$A/(\lambda_1N_0)$},
                        legend pos=north east,
                        grid=both,
                        title={Logarithmisch: $\lambda_2=5\lambda_1$}
                    ]
                    \addplot[domain=0:5,samples=200,thick] {exp(-x)};
                    \addlegendentry{$A_1$}
                    \addplot[domain=0.001:5,samples=200,thick,dashed] {(5/(5-1))*(exp(-x)-exp(-5*x))};
                    \addlegendentry{$A_2$}
                    \end{semilogyaxis}
                \end{tikzpicture}
                \caption{Aktivitäten für $\lambda_2=5\lambda_1$ mit logarithmischer $y$-Achse.}
            \end{figure}

            \begin{notebox}
                Tafelinterpretation:
                \begin{itemize}
                    \item $A_1$ fällt immer exponentiell ab, weil $X_1$ nur zerfällt.
                    \item $A_2$ startet bei $0$, weil am Anfang kein $X_2$ vorhanden ist.
                    \item Danach steigt $A_2$ zuerst an, weil $X_2$ aus $X_1$ produziert wird.
                    \item Später fällt $A_2$ wieder ab, weil die Produktion aus $X_1$ schwächer wird und $X_2$ selbst weiter zerfällt.
                    \item Für $\lambda_2=5\lambda_1$ zerfällt die Tochter schnell; deshalb erreicht $A_2$ rasch ein Maximum.
                    \item Für $\lambda_2=\lambda_1/5$ zerfällt die Tochter langsam; deshalb baut sich $A_2$ langsamer auf und bleibt länger relevant.
                \end{itemize}
            \end{notebox}

    \end{enumerate}

\end{document}
